% ======================================================================
\section{Setting the stage}

% ======================================================================
\subsection{Preprocessor}

% ----------------------------------------------------------------------
\begin{frame}[fragile]{What is it?}

  \begin{itemize}
  \item Commands that (conditionnally) transform text, {\em prior to compilation}
    \begin{itemize}
    \item \lstinline{#include}, \lstinline{#if}, \lstinline{#else}, \lstinline{#elif}, \lstinline{#endif}, \lstinline{#define}, \lstinline{#ifdef}, \lstinline{#ifndef}.
    \end{itemize}
  \item Two output facilities, one to trigger failure:
    \begin{itemize}
    \item \lstinline{#warning "You have been warned"}: compilation continues after printing
    \item \lstinline{#error "teH iMPosSiBLe HaPpEneD!"}: compilation fails after printing
    \end{itemize}
  \item Two operators (your life just changed):
    \begin{itemize}
    \item \lstinline{# symbol}

      will create {\em constant string} \lstinline{"symbol"} (symbol ``stringification'')
    \item \lstinline{foo ## bar}

      will create {\em symbol} \lstinline{foobar} (symbol concatenation)
    \item One minor glitch: better used encapsulated, more later.
    \end{itemize}
  \item What is actually compiled is the output of the preprocessor.
    \begin{itemize}
    \item What about conditional copy-pastes that end up writing code for you?
    \end{itemize}
  \end{itemize}

\end{frame}

% % ----------------------------------------------------------------------
\begin{frame}[fragile]{Example}

  Crudeness of the preprocessor commands makes debugging them a little tedious:

  \begin{lstlisting}
    #include <stdint.h> /* Integers and their bounds */

    #if UINTPTR_MAX == UINT64_MAX
    #warning "64-bit pointers. Now that's a machine!"

    #elif UINTPTR_MAX == UINT32_MAX
    #warning "32-bit pointers. I like little gadgets too!"

    #else
    #error "I will refuse to compile on alien hardware"
    #endif
  \end{lstlisting}

  There's no \lstinline{printf()}-like facility nor any C construct!

  Except for (possibly parenthesized) logic/arithmetic testing with \lstinline{#if} and \lstinline{#elif}.

\end{frame}

% ----------------------------------------------------------------------
\begin{frame}[fragile]{Symbols}
  
  \begin{itemize}
  \item Preprocessor symbols may or may not be defined.
  \item When they are defined, they may or may not have a value.
    \begin{itemize}
    \item And when they have a value, it may or may not be a {\em macro}.
    \end{itemize}
  \end{itemize}

  This will eventually let us achieve C code abstraction w.r.t.
  \begin{itemize}
  \item Machine-dependent features (pointer size, SIMD extensions, {\em etc.})
    \begin{itemize}
    \item Isolate code for bleeding-edge hardware, preserve safe fallback for production.
    \end{itemize}
  \item Lower-level implementation choices:
    \begin{itemize}
    \item \lstinline{set.h} $\rightarrow$ \lstinline{map.h} $\rightarrow$ \{ \lstinline{hashtable.h}, \lstinline{bst.h}, \lstinline{list.h} \}
    \item \lstinline{set.h} will essentially build on \lstinline{put()}/\lstinline{get()} from \lstinline{map.h} to implement union, membership testing, symmetric difference, {\em etc.}

      Any other structure providing \lstinline{put()}/\lstinline{get()} could be used.
    \item \lstinline{map.h} will factor out all possible code so the core algorithmic differences are isolated in \lstinline{list.h}, \lstinline{bst.h} and \lstinline{hashtable.h}.
    \end{itemize}
  \end{itemize}

\end{frame}

% ----------------------------------------------------------------------
\begin{frame}[fragile]{Preprocessor caveats (side-effects)}

  Suppose we have:
  \begin{lstlisting}
    #define gen_abs(x) ((x)<0?-(x):(x))
  \end{lstlisting}

  If you use it with side-effects, a kitten will die:
  \begin{lstlisting}
    int a = -3;

    int gabs = gen_abs(a++);

    /* Expands to:   int gabs = ((a++)<0?-(a++):(a++));
    
       => The value of a is now -1, yet -2 was likely expected.
    */
  \end{lstlisting}

\end{frame}

% ----------------------------------------------------------------------
\begin{frame}[fragile]{Preprocessor caveats (parentheses)}

  Suppose we have:
  \begin{lstlisting}
    #define get_idx(p) p->idx
  \end{lstlisting}

  If you use it with pointer arithmetics, a kitten will die:
  \begin{lstlisting}
    struct foo_t { int idx; };
    
    struct foo_t *arr = calloc(42, sizeof *arr);
    /* Fill arr as expected */
    int id = get_idx(arr+2);

    /* Expands to:   int id = arr+2->idx;
    
       => Syntax error.
    */
  \end{lstlisting}

\end{frame}

% ----------------------------------------------------------------------
\begin{frame}[fragile]{Preprocessor caveats (do and don't)}


  \begin{enumerate}
  \item {\em Never} use a macro with side-effects in the argument(s);
  \item Put parentheses around every parameter occurrence:
  \begin{lstlisting}
    #define get_idx(p) (p)->idx
  \end{lstlisting}
  And around the macro body itself if it's to be used inside other expressions:
  \begin{lstlisting}
    #define list_is_empty(list) ((list) == LIST_END_MARKER)
  \end{lstlisting}
  \end{enumerate}

  {\em Incipit trag\oe dia.}
  
\end{frame}

% ----------------------------------------------------------------------
\begin{frame}[fragile]{Common pattern: Arch.-dependent constants}

  Suppose the start of some \lstinline{src/list.c}:
  
  \begin{lstlisting}
    #include <stdint.h> /* Integers and their bounds */
    /* We may have other plans for NULL */
    #if   UINTPTR_MAX == UINT64_MAX
    /* Define a symbol with a constant value */
    #define LIST_END_MARKER ((void*)0xffffffffffffffffULL)
    #elif UINTPTR_MAX == UINT32_MAX
    #define LIST_END_MARKER ((void*)0xffffffffUL)
    #else
    #error "Unsupported architecture"
    #endif

    /* Define a *macro* symbol, works on 32/64-bit machines */
    #define LIST_END(list) ((list) == LIST_END_MARKER)
  \end{lstlisting}
  
\end{frame}

% ----------------------------------------------------------------------
\begin{frame}[fragile]{Common pattern: Arch.-dependent constants}

  Our \lstinline{src/list.c} may continue like:
  
  \begin{lstlisting}
    #include <stddef.h> /* size_t */

    size_t list_length(list_t l) {
      size_t len = 0;

      for ( ; !LIST_END(l) ; l = l->next)
        len++;

      return len;
    }
    
    /* Plenty of code to finish implementing lists */

    #undef LIST_END         /* Clean locally-defined symbols */
    #undef LIST_END_MARKER
  \end{lstlisting}

\end{frame}

% ----------------------------------------------------------------------
\begin{frame}[fragile]{Common pattern: Arch.-dependent constants (cont'd)}

  What the compiler will automatically get for compilation on a 32-bit machine:
  
  \begin{lstlisting}
    size_t list_length(list_t l) {
      size_t len = 0;

      for ( ; !((l)==(void*)0xffffffffUL) ; l = l->next)
        len++;

      return len;
    }
  \end{lstlisting}

  Use \lstinline{-E} compiler flag to get the preprocessor output.

  {\em Warning!} This is very verbose (all the \lstinline{#include}'d contents is there).

  You have to anticipate for the correct expansion of macros and symbol values!
  
\end{frame}
  
% ----------------------------------------------------------------------
\begin{frame}[fragile]{Common pattern: Conditional inclusion}

  Suppose a higher-level \lstinline{include/structure.h}:
  
  \begin{lstlisting}
    #if defined(TREE_STRUCT)

    #include <structure/tree.h>

    #elif defined(LIST_STRUCT)

    #include <structure/list.h>

    #else
    #error "Unsupported/unspecified data structure"
    #endif /* TREE_STRUCT, LIST_STRUCT */

    /* Common code for all data structures */
  \end{lstlisting}

\end{frame}

% ----------------------------------------------------------------------
\begin{frame}[fragile]{Common pattern: Conditional inclusion}

  Now in some \lstinline{code.c}, we'd certainly expect to be able to do:
  
  \begin{lstlisting}
    #define LIST_STRUCT
    #include <structure.h>
    /* Will eventually include <structure/list.h> */
  \end{lstlisting}

  You may want to see \lstinline{LIST_STRUCT} as a {\em parameter} of the {\em template} \lstinline{structure.h}.

  The idea is:
  \begin{itemize}
  \item to let the user specifiy an option (here: a data structure backend);
  \item to factor out the maximum amount of code among options;
  \item to have the preprocessor write the actual code for us.
  \end{itemize}

  This is the basic paradigm of metaprogramming: programming in program space.
  \begin{itemize}
  \item Scheme, Rust (and many others!) feature much more powerful such capabilities.
  \item The C preprocessor \href{https://github.com/pfultz2/Cloak/wiki/Is-the-C-preprocessor-Turing-complete\%3F}{has internal limits} (and arguably a primitive syntax!)
  \end{itemize}
  
\end{frame}

% ----------------------------------------------------------------------
\begin{frame}[fragile]{Common pattern: Templating (motivation)}

  Suppose you want a map to get the histogram of words (mutations) in a text (\textsc{dna}).
  \begin{itemize}
  \item The map keys are the (unique) words, their value is the word count in the text.
  \end{itemize}

  We want a {\em generic} map template that could be {\em instantiated} more or less like:
  \begin{lstlisting}
    #define key_t        char*
    #define key_hash     word_hash
    #define key_compare  wordptr_compare
    #define value_t      size_t
    #define tag          histogram
    #include <structure/map.h>
  \end{lstlisting}

  \begin{itemize}
  \item Expected result: After line 6, we have the \lstinline{histogram_t} type and we can use \lstinline{histogram()} (the {\em constructor}), \lstinline{histogram_get()}, \lstinline{histogram_put()}, {\em etc.}
  
  \item Encapsulating concatenation is mandatory due to the preprocessor internals:
  \begin{lstlisting}
    #define CONCAT(foo,bar) foo ## bar
    #define CAT(foo,bar)    CONCAT(foo,bar)
  \end{lstlisting}
\end{itemize}

\end{frame}

% ----------------------------------------------------------------------
\begin{frame}[fragile]{Common pattern: Templating}

  Now in \lstinline{map.h}, we first define ({\em prologue}) and then clean ({\em epilogue}) local symbols:

  \begin{lstlisting}
    /* Generic type name: T (map of any key/value types) */
    #define T         CAT(tag,_t)   /* Symbolic prologue */
    #define map_put   CAT(tag,_put)

    typedef struct T *T;     /* Generic type declaration */

    int map_put(T map, key_t key, value_t value) {
      /* Insert/update key/value in map */
    }

    #undef map_put                  /* Symbolic epilogue */
    #undef T
  \end{lstlisting}

  \begin{itemize}
  \item Except for the preprocessor markup, this {\em generic} code looks an awful like plain C!
  \end{itemize}

\end{frame}

% ----------------------------------------------------------------------
\begin{frame}[fragile]{Common pattern: Templating + polymorphism}

  We better acknowledge our map may only be used as a set (no \lstinline{value_t} would exist).

  Lines 7--9 of the previous slide may then become:

  \begin{lstlisting}
    #ifdef value_t
    int map_put(T map, key_t key, value_t value) {
      /* Insert or update key/value mapping in map    */
    #else
    int map_put(T map, key_t key) {
      /* Try to insert key (element) in map (set)     */
    #endif

    /* Code that only relies on the key (most of it!) */

    #ifdef value_t
    /* Do the only difference due to having a value_t */
    #endif

    }
    
  \end{lstlisting}
  
\end{frame}

% ----------------------------------------------------------------------
\begin{frame}[fragile]{Common pattern: Templating + polymorphism (cont'd)}

  Obviously, polymorphism also extends to the definition of C \lstinline{struct}'s:

    \begin{lstlisting}
      #define mapping_t CAT(tag,_mapping_t)

      struct {
        key_t   key;
      #ifdef value_t
        value_t value;
      #endif
      } mapping_t;
      
      /* Don't forget to #undef mapping_t at the end! */
      #ifdef value_t
      #undef tag
      #endif /* If coming from set.h, will be #undef'd there */
  \end{lstlisting}

  \begin{itemize}
  \item Here, no definition for \lstinline{value_t} means using the map as a set!
  \end{itemize}

\end{frame}

% ----------------------------------------------------------------------
\begin{frame}[fragile]{Common pattern: Templating (true code reuse)}

  Now suppose \lstinline{set.h} is to be built from any structure providing \lstinline{put()}/\lstinline{get()}:

  \begin{lstlisting}
    #define T         CAT(tag,_t)
    #define set_put   CAT(tag,_put)

    #include <structure/map.h> /* Get map_put(), no value_t */
    /* Or use conditional inclusion for even more backends! */

    /* Since tag is the same (it is inherited in map.h), both
       - set_put resolves to histogram_put in set.h (here) and
       - map_put resolves to histogram_put in map.h => Profit!
       Hence, only set ops remain to be described:       
     */

    #define set_union CAT(tag,_union)
    
    T  set_union(T set_A, T set_B) {
    }
  \end{lstlisting}
  
\end{frame}

% ----------------------------------------------------------------------
\begin{frame}[fragile]{Common pattern: Packaging}

  \lstinline{static} functions are {\em not visible} outside their file (their {\em Translation Unit}).
  \begin{itemize}
  \item Perfect for small, local functions of a greater, public code.
  \end{itemize}

  A function that lacks the \lstinline{static} qualifier is visible from outside.
  \begin{itemize}
  \item Library code (we need {\em access} to the \lstinline{cos}/\lstinline{exp}/{\em etc.} symbols from the math library.)
  \end{itemize}

  Forward prototypes declaration is either:
  \begin{itemize}
  \item sensible (\lstinline{static}):

    Some of these functions may be mutually recursive. We've got that covered.
  \item or mandatory (library):

    Like in a \lstinline{.h} header file of the kind you were used to before.
  \end{itemize}

\end{frame}

% ----------------------------------------------------------------------
\begin{frame}[fragile]{Common pattern: Packaging (cont'd)}

  Our single \lstinline{map.h} must take provision for both use cases (polymorphism removed):

  \begin{lstlisting}
    #ifndef ds_interface
    #define ds_interface /* Default: visible functions    */
    #endif
    
    /* First part: forward prototype + type declarations. */
    ds_interface int map_put(T map, key_t key, value_t value);

    /* Second part, only if we actually need the code:    */
    #if !defined(export_interface)
    ds_interface int map_put(T map, key_t key, value_t value) {
      /* Blah. */
    }
    #endif

    #undef ds_interface
  \end{lstlisting}
  
\end{frame}

% ----------------------------------------------------------------------
\begin{frame}[fragile]{Common pattern: Packaging (now unboxing!)}

  Going back to our example, actual map code is only needed in \lstinline{histogram.c}:
  \begin{lstlisting}
    #include <stddef.h>  /* size_t */
    #define key_t   char*
    #define value_t size_t
    #define tag     histogram
    #include <structure/map.h>

    double histogram_IQR(histogram_t histogram) {
      /* Blah. */
    }
  \end{lstlisting}

  \begin{itemize}
  \item Observe that the map type gets defined there, so we do {\em not} include \lstinline{histogram.h}.
  \end{itemize}

\end{frame}

% ----------------------------------------------------------------------
\begin{frame}[fragile]{Common pattern: Include guard}
  Now for \lstinline{histogram.h}, assuming at least one map function must be public:
  \begin{lstlisting}
    #ifndef HISTOGRAM_H
    #define HISTOGRAM_H

    #include <stddef.h>  /* size_t */
    #define key_t   char*
    #define value_t size_t
    #define tag     histogram
    #define export_interface /* No code needed this time */
    #include <structure/map.h>

    /* Additional functions: */
    double histogram_IQR(histogram_t histogram);
    
    #endif /* HISTOGRAM_H guard */
  \end{lstlisting}

  The guard prevents redefining the symbols (which is illegal) upon multiple inclusions.

\end{frame}

% ----------------------------------------------------------------------
\begin{frame}[fragile]{Common pattern: Local inlining}

  A \lstinline{static inline}-qualified function {\em may} be inlined by the compiler:

  \begin{itemize}
  \item If its body is worth copy-pasting into the caller, we save the cost of function call.
  \item Inlined code is then like macros, yet at the compiler level (with type-checking!)
  \end{itemize}

  Extremely useful to get the fastest maps:

  \begin{lstlisting}
    static inline int wordptr_compare(const void *_left, const void *_right) {
      return strcmp(*(char**)_left, *(void**)_right);
    }
    
    static inline size_t word_hash(char *w) {
      size_t h = 1;
      for (size_t a = 31 ; *w ; w++ , a *= 31) h += a**w;
      return h;
    }
  \end{lstlisting}
  
\end{frame}

% ----------------------------------------------------------------------
\begin{frame}[fragile]{Common pattern: Local inlining (cont'd)}

  Some of our map functions {\em themselves} may be small.
  \begin{itemize}
  \item For benchmarking, we therefore take advantage of the \lstinline{ds_interface} option.
  \end{itemize}

  After \lstinline{wordptr_compare} and \lstinline{word_hash}, finish instantiating the fastest possible code:
  \begin{lstlisting}
    #define key_t         char*
    #define value_t       size_t
    #define key_compare   wordptr_compare
    #define key_hash      word_hash
    #define ds_interface  static inline
    #include <structure/map.h>
  \end{lstlisting}

  If our comparison and/or hash functions are small enough, then we'll have fast code.

  Like, {\em really} fast.

  Once we've tuned one or two more things!
  
\end{frame}

% ----------------------------------------------------------------------
\begin{frame}[fragile]{Preprocessor recap}

  Combined, the previous techniques allow for:
  \begin{enumerate}
  \item Type-independent (generic) programming;
  \item Generation of the fastest possible code for each actual types;
  \item True reuse and seamless extension of generic code;
  \item Arbitrarily complex data structure implementations;
  \item Sensible user interface.
  \end{enumerate}

  To get them right at the same time for some real code is tedious (but not really hard).

  So we give the 1200 lines of \lstinline{map.h}.
  \begin{itemize}
  \item You now have all the cards to understand them.
  \end{itemize}

  Of course, map tests and benchmarks have all been templated too.
  
\end{frame}

% ----------------------------------------------------------------------
\begin{frame}[fragile]{True preprocessor madness}

  What precedes is purely ``standard practice'' in advanced (and embedded) C.

  \begin{itemize}
  \item More advanced use of the preprocessor may be leveraged for \href{https://github.com/Hirrolot/awesome-c-preprocessor}{even more generality}.
  
  \item Arguably one of the most extreme pushes along these lines is the \href{https://github.com/P-p-H-d/mlib}{\lstinline{M*LIB}} library.
    \begin{itemize}
    \item Even programs get abstracted;
    \item Many (much) more involved techniques inside.
    \end{itemize}
  \end{itemize}

\end{frame}

% ======================================================================
\subsection{Debugging}

% ----------------------------------------------------------------------
\begin{frame}[fragile]{Separation of concerns}

  \begin{itemize}
  \item A search procedure in a heap that ends up exploring a negative backend array index has a {\em broken invariant} (a proper bug);
  \item A file that cannot be found because the user misspelled the name is treated by {\em explicit error handling}.
  \end{itemize}

  Bugs must be detected as early as possible and cause the program to abort.
  \begin{itemize}
  \item Encompasses failed pre-/post-conditions ({\em incl.} failed allocations).
  \end{itemize}

  Errors might be recoverable, should have a dedicated interface.
  
\end{frame}


% ----------------------------------------------------------------------
\begin{frame}[fragile]{Standard debugging}


  Debugging is essentially about actually enforcing your code {\em invariants}.

  \begin{itemize}
  \item By default, code is compiled in debug mode
  \item Release mode: \lstinline{#define NDEBUG} (well-known preprocessor switch)
  \end{itemize}

  Code that is explicitly executed when not in release mode:
  \begin{itemize}
  \item \lstinline{assert(invariant);} (see \lstinline{man 3 assert})
    \begin{itemize}
    \item More useful variant: \lstinline{assert(invariant && "Bug reason")}
    \end{itemize}
  \item With explicit preprocessor markup:
    \begin{lstlisting}
      #ifndef NDEBUG
      fprintf(stderr, "Searching at index: %d\n", idx);
      #endif
    \end{lstlisting}
  \end{itemize}

\end{frame}


% ----------------------------------------------------------------------
\begin{frame}[fragile]{Powerful debugging}

  \begin{itemize}
  \item By default, you want all your invariants be checked;
  \item But you also want to turn on/off extra debugging for the specific module \lstinline{foo}.
  \end{itemize}

  Use a \lstinline{DEBUG_FOO} preprocessor switch symbol:
  \begin{itemize}
  \item Un/comment \lstinline{DEBUG_FOO} definition to de/activate extra debugging code.
  \item Plenty of examples below.
  \end{itemize}

\end{frame}


% ----------------------------------------------------------------------
\begin{frame}[fragile]{Powerful debugging (example)}

  Use builtin \lstinline{__func__} and \lstinline{__LINE__} symbols.
  
  Then in some function belonging to \lstinline{src/foo.c}:

  \begin{lstlisting}
    #include <foo.h>
    // #define DEBUG_FOO // Extra debug switch

    void foo_bar(foo_t foo) {

    #if !defined(NDEBUG) && defined(DEBUG_FOO)
      fprintf(stderr, "%s:%d: Arg foo: %p\n", __func__, __LINE__, foo);
    #endif
      /* Check pre-condition independently of DEBUG_FOO */
      assert(foo!=NULL && "NULL argument for foo_bar");

      /* Actual code for foo_bar */
    }
  \end{lstlisting}
  
\end{frame}


% ======================================================================
\subsection{Error handling}

% ----------------------------------------------------------------------
\begin{frame}[fragile]{Standard error handling}

  \begin{lstlisting}
    #include <stdio.h>  /* fopen()    */
    #include <errno.h>  /* errno      */
    #include <string.h> /* strerror() */

    FILE *fun(char *filename) {
      FILE *fp;

      fp = fopen(filename,"rt");
      if (NULL==fp) {
        fprintf(stderr, "%s:%d: fopen(%s): %s.\n", __func__, __LINE__, filename, strerror(errno));
        return NULL;
        /* Caller only knows about an error, not which one */
      }

      return fp;
    }
  \end{lstlisting}
  
\end{frame}


% ----------------------------------------------------------------------
\begin{frame}[fragile]{Standard error handling (discussion)}

  \begin{itemize}
  \item Nice, system-wide interface
  \item But not all functions will modify \lstinline{errno}
  \item Errors ({\em and their cause!}) do not propagate easily
  \item Applicative code needs a more streamlined abstraction
  \end{itemize}

  General solution: unified type for value {\em or} error (and cause)
  
  \begin{itemize}
  \item Kernel-land trick: use reserved address range to encode error codes as pointers
  \item Below: user-land proposal (portable), as in Rust/OCaml/{\em etc.}
  \end{itemize}
  
\end{frame}


% ----------------------------------------------------------------------
\begin{frame}[fragile]{Powerful error handling}

  Applicative code essentially manipulates pointers:
  \begin{lstlisting}
    typedef struct {
      int  error;
      void *data;
    } rval_t;
  \end{lstlisting}

  \begin{itemize}
  \item \lstinline{error==0} indicates success, with (possibly unused) actual return value in \lstinline{data};
  \item \lstinline{error>0} is an error code, with unused \lstinline{data} (reuse/extend \lstinline{errno} error codes!)
  \end{itemize}
  
  \begin{lstlisting}
    rval_t rval(int error, void *data) {
      return (rval_t){.error=error, .data=data};
    }
    #define OK(data)      rval(0,data)
    #define ERR(errno)    rval(errno,NULL)
    #define IS_ERR(rv)  ((rv).error > 0)
    #define RVAL(rv)     (rv).data
    #define RCODE(rv)    (rv).error
  \end{lstlisting}

\end{frame}


% ----------------------------------------------------------------------
\begin{frame}[fragile]{Powerful error handling, with optional return flow tracing}

  \begin{lstlisting}
    // #define DEBUG_RETFLOW /* Uncomment to print returns */
    typedef struct {
      int  error;
      void *data;
    } rval_t;
    
    #if !defined(NDEBUG) && defined(DEBUG_RETFLOW)
    rval_t rval(int error, void *data, char *func, int line);
    #define OK(data)      rval(0,data,__func__, __LINE__)
    #define ERR(errno)    rval(errno,NULL,__func__,__LINE__)
    #else
    rval_t rval(int error, void *data);
    #define OK(data)      rval(0,data)
    #define ERR(errno)    rval(errno,NULL)
    #endif
    /* Predicate IS_ERR and accessor RVAL keep the same */
  \end{lstlisting}

\end{frame}


% ----------------------------------------------------------------------
\begin{frame}[fragile]{Powerful error handling, with optional return flow tracing (cont'd)}

  \begin{lstlisting}
    rval_t rval(int error, void *data
    #ifndef NDEBUG
                , char *func, int line
    #endif
                ) {
      
    #if !defined(NDEBUG) && defined(DEBUG_RETFLOW)
    #include <stdio.h>
      fprintf(stderr, "[%16s:%d]  return ", func, line);
      if (error>0) fprintf(stderr, "ERROR(%d) ", error);
      fprintf(stderr, "\n");
    #endif

      return (rval_t){.error=error, .data=data};
    }
  \end{lstlisting}

\end{frame}


% ======================================================================
\subsection{Memory management}

% ----------------------------------------------------------------------
\begin{frame}[fragile]{Standard memory management}

  \begin{lstlisting}
    #include <stdlib.h>
    
    int main(void) {
      void *ptr = malloc(42);
      
      free( ptr );
      
      printf("ptr=%p\n", ptr);
      
      exit(EXIT_SUCCESS);
    }
  \end{lstlisting}
  
  \begin{itemize}
  \item Fine-grained allocations only
  \item Should check for \lstinline{ptr!=NULL}
  \item Not easy to trace, debug, nor extend
  \item \lstinline{ptr} keeps its value after the call to \lstinline{free()}
  \end{itemize}
  
\end{frame}


% ----------------------------------------------------------------------
\begin{frame}[fragile]{Powerful memory management: allocators}

  \begin{itemize}
  \item Store allocator info before block.
    \begin{itemize}
    \item Default \lstinline{NULL} allocator maps to standard \lstinline{malloc()}/\lstinline{free()} allocator.
    \end{itemize}
  \item Enforce \lstinline{NULL} after deallocation.
  \end{itemize}

  \begin{lstlisting}
    // #define DEBUG_MEM /* Uncomment to debug mem de/allocs */
    
    typedef struct {
      void   *state; /* Allocator internal data */
      void *(*allocate)(size_t size, void *state);
      void  (*deallocate)(void **object_ptr, void *state);
    } allocator_t;
  \end{lstlisting}

  Ultimately, make any object {\em propagate} its allocation policy:
  \begin{itemize}
  \item Allocate rare objects with \lstinline{NULL} allocator;
  \item Draw tree nodes from {\em e.g.}, an arena allocator
    \begin{itemize}
    \item Coarse-grained memory management useful for writing parsers, games, or GUI;
    \item Key idea: delete all objects (AST subtrees, GUI widgets, {\em etc.}) in arena {\em in $O(1)$ time}.
    \end{itemize}
  \end{itemize}
  
\end{frame}


% ----------------------------------------------------------------------
\begin{frame}[fragile]{Powerful memory management: allocators (interface)}

  \begin{lstlisting}
    #if !defined(NDEBUG) && defined(DEBUG_MEM)
    #include <stdio.h>
    void   *alloc__(size_t size, allocator_t *allocator, char *fun, int line);
    #define alloc(size,al) alloc__(size,al,__func__, __LINE__)
    void    dealloc__(void **objp, char *fun, int line);
    #define dealloc(objp)  dealloc__(objp,__func__, __LINE__)
    #else
    void   *alloc__(size_t size, allocator_t *allocator);
    #define alloc(size,al) alloc__(size,al)
    void    dealloc__(void **objp);
    #define dealloc(objp)  dealloc__(objp)
    #endif /* Augment alloc()/dealloc() with shortcuts: */    
    #define block(sz) alloc(sz, NULL)
    #define new(obj)  block(sizeof(*obj))
    #define like(obj) alloc(sizeof(*obj),((allocator_t*)obj)-1)
    #define del(objp) dealloc((void**)(objp))
  \end{lstlisting}
  
\end{frame}


% ----------------------------------------------------------------------
\begin{frame}[fragile]{Powerful memory management: allocators (alloc)}

  \begin{lstlisting}
    void   *alloc__(size_t size, allocator_t *al
    #if !defined(NDEBUG) && defined(DEBUG_MEM)
                    , char *func, int line
    #endif
                   ) { /* Add space for pointer to allocator */
      allocator_t *ptr = al ?
                    al->allocate(size+sizeof(ptr),al->state) : calloc( 1, size+sizeof(ptr) );
      assert( ptr && "Alloc failed!" ); /* Univ. precond. */
      *ptr = al; /* Propagate allocator info */
    
    #if !defined(NDEBUG) && defined(DEBUG_MEM)
      fprintf(stderr, "[%16s:%d]  alloc %zu bytes at %p\n", func, line, size, ptr+1);
    #endif
      return ptr+1; /* Return address *after* allocator info */
    }
  \end{lstlisting}
  
\end{frame}


% ----------------------------------------------------------------------
\begin{frame}[fragile]{Powerful memory management: allocators (dealloc)}

  \begin{lstlisting}
    void   *dealloc__(void **objp
    #if !defined(NDEBUG) && defined(DEBUG_MEM)
                      , char *func, int line
    #endif
                     ) { /* Recover allocator info from objp */
      if (objp && *objp) {
        allocator_t **alp = ((allocator_t**)*objp)-1;
    #if !defined(NDEBUG) && defined(DEBUG_MEM)
        fprintf(stderr, "[%16s:%d]  dealloc at %p\n", *object_ptr, fun, line);
    #endif
        if (*alp)
          (*alp)->deallocate(alp, (*alp)->state);
        else free(alp);
        *objp = NULL; /* No silver bullet, yet sane overall */
      }
    }
  \end{lstlisting}
  
\end{frame}


% ----------------------------------------------------------------------
\begin{frame}[fragile]{Powerful memory management: extensions}

  \begin{enumerate}
  \item In \lstinline{de/alloc__()}, add global number and total size of live memory blocks
    \begin{itemize}
    \item Monitor program memory usage, report statistics
    \end{itemize}
  \item In \lstinline{de/alloc__()}, add global {\em set} of live memory blocks
    \begin{itemize}
    \item Possibly with function/line info;
    \item Program can reflect on the memory it uses;
    \item Main application: interned (constant, unique) strings (Python VM);
    \item Cosmetic application: no leaks on exit (register cleanup handler with \lstinline{atexit(3)})
    \end{itemize}
  \item Specialized, fast allocators for objects of known, constant size
    \begin{itemize}
    \item List links, tree nodes, {\em etc.} (we'll do that for maps!)
    \end{itemize}
  \end{enumerate}
  
  Overall goal: minimize number of calls to the default system allocator
  \begin{itemize}
  \item Ultimately, any allocator will have to use it anyway!
  \item List of reusable blocks, {\em etc.}
  \end{itemize}

  At least, the infrastructure is ready for all of this.

\end{frame}


% ----------------------------------------------------------------------
\begin{frame}[fragile]{Recap}

  All of the following also comes with switchable, extra debugging capabilities:
  \begin{itemize}
  \item Algorithms must check their invariants with \lstinline{assert()};
  \item Procedural code uses \lstinline{rval_t} for control management;
  \item Either you get a block from specified allocator, or program aborts.
  \end{itemize}

  Key techniques used:
  \begin{itemize}
  \item Ternary \lstinline{if-then-else} operator: \lstinline{test ? consequence : alternative}
  \item Build on top of standards (\lstinline{NDEBUG}, \lstinline{errno}), seek retro-compatibility
  \item Preprocessor (combined) switches, predefined \lstinline{__func__} and  \lstinline{__LINE__} macros
    \begin{itemize}
    \item \lstinline{__FILE__} also available for more fancy bloat;
    \item Replace the \lstinline{fprintf(stderr, ...)} by the colorful \lstinline{INFO()}, \lstinline{WARNING()}, {\em etc.} macros in \lstinline{common/logging.h};
    \end{itemize}
  \item Guarded access to pointers: \lstinline{if (foo) foo->bar = 0;}
  \item Block metadata: adding allocator pointer {\em before} returned address (more below).
  \end{itemize}
  
\end{frame}


% ----------------------------------------------------------------------
\begin{frame}[fragile]{Example (load file contents into memory)}

  \begin{lstlisting}
    rval_t file_contents(char *filename, size_t *size) {
      FILE       *fp;
      char    *bytes;
      struct stat sb;
      assert(filename && "Missing filename" );
      if (-1==stat(filename, &sb)) return ERR(errno);
      if (!(fp = fopen(filename, "rt"))) return ERR(errno);
      bytes = block(sb.st_size);
      for ( size_t r = 0, rr = 0 ; r < sb.st_size ; r += rr ) {
        rr = fread(bytes+r, sb.st_size-r, 1, fp);
        if (ferror(fp)) { del(&bytes); fclose(fp);
          return ERR(ENODATA);
        }
      }
      fclose(fp); if (size) *size = sb.st_size;
      return OK(bytes);
    }
  \end{lstlisting}
  
\end{frame}


% ----------------------------------------------------------------------
\begin{frame}[fragile]{Example (using it)}

  \begin{lstlisting}
    int main(int argc, char *argv[]) {
      size_t   length;
      char  *filename = argc>1 ? argv[1] : NULL;
      rval_t contents = file_contents( filename, &length );

      if (IS_ERR(contents)) {
        fprintf(stderr, "%s: FAILED (error %d).\n", filename, RCODE(contents));
        exit(RCODE(contents));
      }
      
      printf("%s: loaded %zu bytes.\n", filename, length);
      del(&RVAL(contents));
      
      exit(EXIT_SUCCESS);
    }
  \end{lstlisting}

\end{frame}


% ----------------------------------------------------------------------
\begin{frame}[fragile]{Example (running it)}

  With full debug trace on error, watch error code propagate to the terminal:

\begin{verbatim}
$ make test/0-core@check
./test/0-core -v
[   file_contents:16]  return ERROR(2) 
-v: FAILED (error 2).
make: *** [.Makefile.inc:233: test/0-core@check] Error 2
$ echo $?
2
\end{verbatim}

  Same on success:

\begin{verbatim}
$ ./test/0-core ./test/0-core.c
[   file_contents:21]  alloc 491 bytes at 0x515000000308
[   file_contents:33]  return 
./test/0-core.c: loaded 491 bytes.
[            main:19]  dealloc at 0x515000000308
\end{verbatim}
  
\end{frame}


